\section*{Data}

To study intra-state conflict patterns in Nigeria we utilize the Armed Conflict Location and Event Data Project (ACLED) initiated by \citet{raleigh:etal:2010}. This dataset records armed conflict and protest events in over 60 developing countries. ACLED's \textit{battles} data is used to generate our measure of conflict where $y_{ij,t} = 1$ indicates that a conflict occurred when actor $i$ attacked actor $j$ at time $t$ ($y_{ij,t} = 0$ if no conflict occurred).  We focus only on armed groups that are engaged in battles for at least 5 years during the 2000-2016 period, which results in a total of 37 armed groups. 

The groups included in our battle-data analysis are politically violent actors as defined by the ACLED codebook. Such actors include rebels, militias, ethnic groups, and civilians.\footnote{For detailed definitions of these actors, see the on-line ACLED Africa data codebook maintained at \url{www.acleddata.com}.} Additionally, we utilized secondary resources to confirm the existence and approximate location of each actor in the data.
%\footnote{Notes and resources about each are included in the Appendix.} 
Overall, the data capture several key types of actors: government armed forces (military and police), insurgent groups (Boko Haram), and ethnic militias (including the Yoruba and Fulani Militias). 

The ACLED data also allows us to explore civilians' role in shaping conflict within the Nigerian system.  Building on the research agenda motivated by macro-level studies \citep{gurr:1970, opp:1988, tarrow:1994, tucker:2007,chenoweth:stephan:2011}, we investigate the link between civilian mobilization and violence between armed actors at the local level by creating a count of the number of protests/riots led by civilians against a given actor at time $t-1$. While there is a robust debate over what causes civilian victimization,\footnote{For example, see \citet{valentino:etal:2004, downes:2006, kalyvas:2006, salehyan:etal:2014, prorok:appel:2014, humphreys:weinstein:2006}.} discussion of the consequences of civilian victimization, particularly as they relate to the potential for future conflict has been more limited \citep{hultman:2007, raleigh:2012}. \citet{berman:matanock:2015} argue that when rebels groups [governments] kill civilians, other civilians are more [less] likely to share information with the government: information allows the government to carry out attacks against insurgents, and the lack of information makes insurgent attacks less likely. \citet{condra:shapiro:2012} find support for this relationship in Iraq.\footnote{\citet{lyall:2009}, however,  a contrary effect and demonstrates that indiscriminate violence against civilians in Chechnya is associated with a lower risk of rebel attacks, possibly because these attacks deprive insurgents of resources. } To account for both these dynamics we create a count of the number of violent actions that an actor committed against civilians at time $t-1$ and nodal covariates to explain both the probability of an actor sending and receiving a conflictual tie. 

Additionally, we include a control for spatial effects for modeling the likelihood of conflict between a particular pair of actors. First, we add a measure of how dispersed a rebel group's activity is across the country. Armed groups whose actions are more dispersed across Nigeria are more likely to come into conflictual contact with others. We include a ``Geographic Spread'' variable as both a sender and receiver covariate.\footnote{For example, to model the spread of group $i$ in period $t$, we calculate this by taking all the events involving group $i$ in the 3 years prior to $t$ and estimate the variance.}  Second, a notable literature has developed explicating the mechanisms through which civil conflict can diffuse between countries.\footnote{For example, see \citet{starr:most:1983,salehyan:2009, salehyan:gleditsch:2007,metternich:etal:2015, braithwaite:2010a, beardsley:2011}.} To account for this effect, we use the ACLED dataset to count the number of conflict events occurring within Nigeria's contiguous neighbors. 

We add three other covariates to the model. First, we create a control for whether both actors are a part of the government. There are only two actors within the 37 that we study in the Nigerian system that would fall under this classification, namely, the Police and Military. We include this control to account for the fact that the probability of an interaction between these two is minimal relative to the other potential dyads in this system.\footnote{An alternative approach would be to include sender and receiver level government indicator variables, our results are similar when using that approach.} Second, we include a binary variable that is one if the following year is an election year and zero otherwise. This is to account for the fact that elections in Nigeria are rarely waged peacefully in the ballot box, but instead are typically followed or preceded by episodes of conflict. Most importantly, to examine the affect that Boko Haram's entrance has had on the level of conflict in this system we include a binary variable that takes on the value of one after the Boko Haram insurgency has begun (in 2009) and zero beforehand.\footnote{We choose this date both because 2009 is considered the beginning of the Boko Haram insurgency in most media reports, despite their founding in 2002,  because 2009 features the first battle including Boko Haram in the ACLED Dataset.}

\section*{Modeling Approach}

To estimate conflict in this system in a way that accounts for interdependencies between actors, we rely on a network based approach that combines the social relations regression model (SRRM) -- for details on this model see \citet{li:loken:2002,hoff:ward:2004,dorff:minhas:2016} -- and the latent factor model developed by \citet{hoff:2009}. Together the SRRM provides a set of additive effects that we can use to capture first and second order dependencies, and the latent factor model provides a set of multiplicative effects that we use to model third order dependencies. We refer to this estimator as the additive and multiplicative effects (AME) model:

\begin{align}
	\begin{aligned}
		y_{ij,t} &= g(\theta_{ij,t}) \\ 
		&\theta_{ij,t} = \bm\beta_{d}^{T} \mathbf{X}_{ij,t} + \bm\beta_{s}^{T} \mathbf{X}_{i,t} + \bm\beta_{r}^{T} \mathbf{X}_{j,t} + e_{ij,t} \\
		&e_{ij,t} = a_{i} + b_{j}  + \epsilon_{ij} + \alpha(\textbf{u}_{i}, \textbf{v}_{j}) \text{  , where } \\
		&\qquad \alpha(\textbf{u}_{i}, \textbf{v}_{j}) = \textbf{u}_{i}^{T} \textbf{D} \textbf{v}_{j} = \sum_{k \in K} d_{k} u_{ik} v_{jk} \\ 
	\label{eqn:ame}
	\end{aligned}
\end{align}

where $y_{ij,t}$ represents whether or not conflict occurred between actor $i$ (the sender) and actor $j$ (the receiver) at time $t$. To model our binary dependent variable we employ a latent variable representation of a probit regression framework, in which we model a latent variable, $\theta_{ij}$, using a set of time varying dyadic ($\bm\beta_{d}^{T} \mathbf{X}_{ij,t}$), sender ($\bm\beta_{s}^{T} \mathbf{X}_{i,t}$), and receiver covariates ($\bm\beta_{r}^{T} \mathbf{X}_{j,t}$). Then we decompose the error using the SRRM and latent factor model. $a_{i}$ and $b_{j}$ represent sender and receiver random effects that we incorporate from the SRRM framework:

\begin{align}
	\begin{aligned}
		\{ (a_{1}, b_{1}), \ldots, (a_{n}, b_{n}) \} &\simiid N(0,\Sigma_{ab}) \\ 
		\{ (\epsilon_{ij}, \epsilon_{ji}) : \; i \neq j\} &\simiid N(0,\Sigma_{\epsilon}), \text{ where } \\
		\Sigma_{ab} = \begin{pmatrix} \sigma_{a}^{2} & \sigma_{ab} \\ \sigma_{ab} & \sigma_{b}^2   \end{pmatrix} \;\;\;\;\; &\Sigma_{\epsilon} = \sigma_{\epsilon}^{2} \begin{pmatrix} 1 & \rho \\ \rho & 1  \end{pmatrix}
	\label{eqn:srm}
	\end{aligned}
\end{align}

The interpretation of these parameters is straightforward. The sender and receiver random effects are modeled jointly from a multivariate normal distribution to account for correlation in how active an actor is in sending and receiving ties. Heterogeneity in the the sender and receiver effects is captured by $\sigma_{a}^{2}$ and $\sigma_{b}^{2}$, respectively, and $\sigma_{ab}$ describes the linear relationship between these two effects (i.e., whether actors who send [receive] a lot of ties also receive [send] a lot of ties). Beyond these first-order dependencies, second-order dependencies are described by $\sigma_{\epsilon}^{2}$ and a within dyad correlation, or reciprocity, parameter $\rho$. 

While the additive effects from the SRRM can deal with first (differing popularity and activity of actors) and second order interdependencies (reciprocity), the multiplicative effects are used to deal with third order dependencies. Specifically, this multiplicative effect allows us to model homophily -- the tendency of actors with similar characteristics to form strong relationships than those with differing characteristics -- and stochastic equivalence, the possibility that two actors $i$ and $j$ will have similar relationships with every other actor in the network. An AME model accounts for these third order effects using the multiplicative term: $\alpha(\textbf{u}_{i}, \textbf{v}_{j})$. This model posits a latent vector of characteristics $\mathbf{u_{i}}$ and $\mathbf{v_{j}}$ for each sender $i$ and receiver $j$. The similarity or dissimilarity of these vectors will then influence the likelihood of activity, and therefore account for these third order interdependencies \citep{minhas:etal:2016:arxiv}.\footnote{Parameter estimation in the AME takes place by sampling from the posterior distribution of the full conditionals. Details on the estimation procedure can be found in the Appendix.} 
